\section{Изпълнителни модели на свързващите компоненти и платформено-специфични клиентски библиотеки}
\label{sec:backend_scenarios}

Настоящата глава изследва не само как една и съща C++ абстракция се материализира в различни целеви системи (свързващи компоненти), но и как се изпълнява през конкретните платформено-специфични клиентски библиотеки. Анализът има две цели. Първо, да покаже, че библиотеката не е ограничена до чисто SQL интерпретация. Второ, да изясни къде асинхронната интеграция може да следва платформено-специфичен неблокиращ модел и къде честният архитектурен избор остава изнасяне към набор от нишки (изнасяне към набор от нишки) върху синхронен драйвер.

След теоретичната рамка от Глава~\ref{sec:theoretical_background} и архитектурното обяснение на асинхронния слой от Глава~\ref{sec:async_execution_architecture}, тук вече не се пита дали различните модели за съхранение са еквивалентни. Вместо това се пита дали една и съща логическа ОРС форма може да бъде адаптирана контролирано към тях, без библиотеката да обещава семантика на изпълнение, която даден компонент не може естествено да изпълнява. Тъкмо в това се проявява стойността на базирания на възможности договор на конектора и на разграничението между универсален асинхронен интерфейс и специфична за компонента реализация.

\subsection{Методика на анализа}

Всеки сценарий следва една и съща последователност: теоретична опора, логическо картографиране от C++ модела, реализация в \code{connector\_trait<DB>}, ограничения и верификация чрез тестове. Това прави сравнението между компонентите възможно. Не се търси твърдение, че всички системи са еквивалентни, а че една и съща логическа ОРС абстракция може да бъде адаптирана контролирано към различни модели за съхранение.

\begin{table}[H]
\centering
\small
\caption{Еднаква аналитична рамка за всеки сценарий на свързващ компонент}
\label{tab:backend_analysis_framework}
\begin{tabularx}{\textwidth}{L{3.3cm}Y}
\toprule
\textbf{Критерий} & \textbf{Какво се оценява} \\
\midrule
Теоретична опора & към кой модел за съхранение от Глава~\ref{sec:theoretical_background} принадлежи свързващият компонент \\
Логическо картографиране & как \code{property}, \code{field}, \code{Rule} и \code{table\_name\_trait} се интерпретират в конкретната система \\
Профил на възможностите & кои ОРС операции са изрично допустими и кои трябва да бъдат отхвърлени \\
Поведение при изпълнение & как се рендират заявки, параметри и резултатни проекции \\
Тестова опора & кои модулни (модулни) и интеграционни тестове доказват договора в хранилището \\
\bottomrule
\end{tabularx}
\end{table}

Тази аналитична рамка е важна, защото позволява сравнението да остане честно. SQL системите не се привилегироват само защото имат по-богат синтаксис, а не-SQL (NoSQL) системите не се оценяват като ``непълни'' само защото отказват операции, които са неподходящи за техния модел на съхранение.

\begin{figure}[H]
\centering
\begin{tikzpicture}[node distance=1.35cm and 1.2cm]
    \node[ormaccent] (ir) {\textbf{Споделено междинно представяне}\\полета, предикати, заместители, проекции};
    \node[ormblock, below left=of ir] (sql) {SQL материализация\\SQLite/PostgreSQL/MySQL};
    \node[ormblock, below=of ir] (doc) {Документна материализация\\MongoDB};
    \node[ormblock, below right=of ir] (graph) {Графова материализация\\Neo4j};
    \node[ormblock, below=of doc, xshift=-3.2cm] (redis) {Материализация по ключ\\Redis};
    \node[ormblock, below=of doc, xshift=3.2cm] (cass) {Материализация по разделяне (разделяне)\\Cassandra};

    \draw[ormline] (ir) -- (sql);
    \draw[ormline] (ir) -- (doc);
    \draw[ormline] (ir) -- (graph);
    \draw[ormline] (ir) -- (redis);
    \draw[ormline] (ir) -- (cass);
\end{tikzpicture}
\caption{Едно споделено междинно представяне и множество специфични за компонентите пътища за материализация}
\label{fig:backend_materialisation_paths}
\end{figure}

Фигура~\ref{fig:backend_materialisation_paths} показва ключовото твърдение на главата: не съществуват множество ОРС библиотеки в рамките на хранилището, а една логическа библиотека с различни пътища за материализация. Това е силният аргумент в полза на архитектурата на конектора --- разнообразието от целеви системи е овладяно чрез общо представяне и изричен договор (изричен договор), а не чрез копиране на програмни интерфейси.

\subsection{Релационен базов сценарий: SQLite}

Този сценарий стъпва върху релационния модел, разгледан в Глава~\ref{sec:theoretical_background}. \code{SQLiteDB} е най-естествената референтна целева система за библиотеката, защото съчетава реална SQL семантика, малък експлоатационен контекст и пряка интеграция чрез програмния интерфейс на \code{sqlite3}. В него \code{property} се материализира като колона, \code{table\_name\_trait} --- като име на таблица, а дърветата \code{Rule} --- като \code{WHERE} изрази.

Конекторът декларира \code{supports\_joins}, \code{supports\_transactions} и \code{supports\_aggregation}, което го прави най-близък до класическото ОРС очакване. Именно върху него се валидират подготвени заявки, логика за свързване на параметри (свързване), \code{find\_one()} сценарии и попълване на резултати (попълване на резултати) към типизирани кортежи. Интеграционните тестове в \path{tests/интеграционни/test_sqlite.cpp} показват не само коректност на \code{SELECT}, \code{INSERT}, \code{UPDATE} и \code{DELETE}, но и поведение при конектор с възможност за транзакции и профил с възможност за съединяване (с възможност за съединяване).

От академична гледна точка SQLite играе ролята на релационна базова линия (базова линия). Ако дадена ОРС форма е проблематична още тук, тя трудно би могла да бъде обобщена към по-екзотични системи. Затова настоящата работа използва SQLite като отправна точка, а не просто като един от многото примери.

\subsection{Релационна преносимост: PostgreSQL и MySQL}

Теоретичната основа и тук е релационният модел, но архитектурният интерес е различен: доколко едно и също междинно представяне остава валидно при различни SQL диалекти и различни правила за свързване на параметри. \code{PostgreSQLDB} и \code{PostgreSQLLiveDB} поддържат \code{supports\_joins}, \code{supports\_transactions} и \code{supports\_aggregation}, а тестовете в \path{tests/модулни/test_postgresql_конектор.cpp} и \path{tests/интеграционни/test_postgresql_live.cpp} показват реално изпълнение срещу жива база данни.

\code{MySQLDB} и \code{MySQLLiveDB} също декларират \code{supports\_joins} и \code{supports\_transactions}, но тук е важно различието в модела за свързване. Докато PostgreSQL може естествено да адресира параметри чрез \code{$1}, \code{$2}, \ldots{}, MySQL разчита на позиционни \code{?} заместители (заместители). Това не нарушава общото междинно представяне, но показва, че логическата преносимост не означава идентична байт по байт материализация в целевата система.

Тъкмо тези сценарии от SQL семейството защитават по-силна теза от ``работи със SQL'': библиотеката демонстрира, че междинното представяне е достатъчно общо за повече от един диалект, а в същото време е достатъчно честно, за да остави нисконивовите различия на конекторния слой.

\subsection{Документен сценарий: MongoDB}

Този сценарий стъпва върху документния модел от теоретичната глава. При MongoDB \code{property} се интерпретира като поле в документ, а дърветата от предикати на заявките се превеждат към BSON-подобни филтри. \code{select} изразите се материализират като проекция (проекция), а таблицата от гледна точка на обектно-релационното съпоставяне се превръща в колекция. Точно тук се вижда практическата стойност на независимото от съхранението междинно представяне: потребителят не описва \code{\$eq}, \code{\$and} и проекция директно, но конекторът може да ги изведе от същата логическа форма.

Важна подробност е, че \code{connector\_trait<MongoDB>} не декларира labelи за възможности в SQL стил. Това не е пропуск, а коректно описание на договора на конектора. Няма претенция, че Mongo целевата система е еквивалентна на съединяване (еквивалентна на съединяване) спрямо релационните конектори. Вместо това той предлага документно-естествено тълкуване на междинното представяне. Модулните и интеграционните тестове в \path{tests/модулни/test_mongodb_конектор.cpp} и \path{tests/интеграционни/test_mongodb_live.cpp} валидират тази адаптация срещу реален драйверен слой.

Архитектурната стойност на този сценарий е, че доказва едно важно ограничение: една библиотека може да работи с множество свързващи компоненти, без да твърди, че всички те поддържат едни и същи операции. Именно този тип ``контролирана частичност'' прави дизайна научно защитим.

\subsection{Сценарий ключ-стойност: Redis}

Този сценарий стъпва върху теорията за ключ-стойност базите данни. При Redis един и същ модел на същността може да бъде материализиран по два основни начина: като \code{SET}/\code{GET} при еднополеви структури или като \code{HSET}/\code{HGETALL} при многополеви структури. \code{table\_name()} участва в генерирането на префикс на ключ, а достъпът е концентриран около равенство по първичен ключ (равенство по първичен ключ).

Тук ясно се вижда, че споделеното междинно представяне не означава идентична изразителност. \code{RedisDB} умишлено не декларира \code{supports\_joins}, \code{supports\_transactions} или \code{supports\_aggregation}. Това ограничение е част от коректния договор, а не слабост на архитектурата. Модулните тестове и тестовете на живо в \path{tests/модулни/test_redis_конектор.cpp} и \path{tests/интеграционни/test_redis_live.cpp} доказват, че библиотеката правилно материализира достъп по ключ, без да обещава неподдържани операции.

Този сценарий е особено ценен за дипломната работа, защото принуждава архитектурата да прави ясно разграничение между \emph{логически модел на данните} и \emph{допустим модел на достъп}. Това разграничение е лесно да се забрави при библиотеките, ориентирани първо към SQL, но тук е изведено като централен принцип.

\subsection{Графов сценарий: Neo4j}

Сценарият за Neo4j стъпва върху графовия модел. Типът същност се материализира като label и набор от свойства, а междинното представяне на заявките се превежда към \code{MATCH}, \code{RETURN} и \code{WHERE} фрагменти. Това е важен архитектурен тест, защото семантиката на графовите заявки е най-отдалечена от класическия SQL модел. Въпреки това логическите елементи --- полета, предикати, филтриране и резултатни проекции --- остават разпознаваеми.

\code{Neo4jDB} и \code{Neo4jLiveDB} декларират \code{supports\_transactions}, но не и възможностите за съединяване/агрегация на SQL слоя. Това е смислено: обхождането на граф не е същото като SQL съединяване, макар на логическо ниво и двете да свързват данни. Модулните и интеграционните тестове в \path{tests/модулни/test_neo4j_конектор.cpp} и \path{tests/интеграционни/test_neo4j_live.cpp} показват, че абстракцията може да адаптира един и същ логически модел и към графова целева система, ако конекторът поеме отговорността за правилната материализация.

От гледна точка на тезата това е една от най-силните демонстрации, че междинното представяне е семантичен, а не синтактичен слой. Ако представянето беше просто недо-SQL, Neo4j сценарият би се разпаднал. Фактът, че е възможна работеща графова материализация, подкрепя централната архитектурна претенция.

\subsection{Сценарий за широки колонни бази: Cassandra}

Този сценарий стъпва върху теорията за широки колонни бази данни (широки колонни). Макар Cassandra да предлага CQL синтаксис, наподобяващ SQL, действителният модел на достъп е подчинен на ключовете за разделяне (ключове за разделяне) и на разпределената организация на данните. Затова библиотеката не може да третира Cassandra като обикновен SQL компонент. Наличието на \code{WHERE} клауза само по себе си не е достатъчно; тя трябва да бъде съвместима с логиката на разделянето.

\code{CassandraDB} и \code{CassandraLiveDB} декларират \code{supports\_transactions}, но не и \code{supports\_joins}. Това е много показателно. От една страна, свързващият компонент продължава да е част от същата ОРС библиотека. От друга страна, договорът категорично отказва операции, които биха били подвеждащи или скъпи в среда на широки колони. Модулните тестове и тестовете на живо в \path{tests/модулни/test_cassandra_конектор.cpp} и \path{tests/интеграционни/test_cassandra_live.cpp} дават конкретна верификация на този по-строг експлоатационен профил (експлоатационен профил).

Именно поради това Cassandra е най-добрият пример, че разширяемостта на библиотеката не е базирана на фалшиво уеднаквяване, а на формално описани граници на допустимите операции.

\begin{figure}[H]
\centering
\begin{tikzpicture}[node distance=1.2cm and 0.95cm]
    \node[ormaccent, minimum width=3.4cm] (title) {Спектър на възможностите};
    \node[ormblock, below left=1.5cm and 2.3cm of title, minimum width=2.8cm] (sqlite) {SQLite\\съединяване + транзакции + агрегация};
    \node[ormblock, below=1.5cm of title, minimum width=3.1cm] (pgsql) {PostgreSQL/MySQL\\съединяване + транзакции + агрегация/транзакции};
    \node[ormblock, below right=1.5cm and 2.3cm of title, minimum width=2.8cm] (mongo) {MongoDB\\без labelи за SQL възможности};
    \node[ormblock, below=of sqlite, minimum width=2.8cm] (redis) {Redis\\само базирано на ключове};
    \node[ormblock, below=of pgsql, minimum width=2.8cm] (neo4j) {Neo4j\\само транзакции};
    \node[ormblock, below=of mongo, minimum width=2.8cm] (cass) {Cassandra\\само транзакции};

    \draw[ormline] (title) -- (sqlite);
    \draw[ormline] (title) -- (pgsql);
    \draw[ormline] (title) -- (mongo);
    \draw[ormline] (sqlite) -- (redis);
    \draw[ormline] (pgsql) -- (neo4j);
    \draw[ormline] (mongo) -- (cass);
\end{tikzpicture}
\caption{Сравнителен профил на възможностите на свързващите компоненти в хранилището}
\label{fig:backend_capability_spectrum}
\end{figure}

Фигура~\ref{fig:backend_capability_spectrum} синтезира едно от най-важните наблюдения в главата: библиотеката не налага еднакъв профил на възможности на всички свързващи компоненти. Точно обратното --- различията се използват като първокласна (първокласна) архитектурна информация.

\subsection{Синтезирана матрица на възможностите}

За да остане сравнението не само качествено, но и таблично проследимо, Фигура~\ref{fig:backend_capability_spectrum} може да бъде сведена до компактна матрица на възможностите. Тя не бива да се чете като универсална истина за самите системи за бази данни, а като описание на конкретната архитектурна позиция на библиотеката спрямо тях в разглежданата ревизия на хранилището.

\begin{table}[H]
\centering
\small
\caption{Профили на възможностите на основните свързващи компоненти според наличните тестове и специализации на конектори}
\label{tab:backend_capability_matrix}
\begin{tabularx}{\textwidth}{L{2.0cm}L{1.45cm}L{1.55cm}L{1.5cm}L{1.35cm}L{1.55cm}Y}
\toprule
\textbf{Компонент} & \textbf{Съединявания} & \shortstack{\textbf{Транзак-}\\\textbf{ции}} & \shortstack{\textbf{Агрега-}\\\textbf{ция}} & \textbf{Обновяване с вмъкване (upsert)} & \shortstack{\textbf{Масово записване}\\\textbf{(Bulk insert)}} & \textbf{Кратък коментар} \\
\midrule
SQLite & налична & налична & налична & \shortstack{ограни-\\чена} & \shortstack{ограни-\\чена} & най-близкият пълен релационен референтен път в хранилището \\
PostgreSQL & налична & налична & налична & \shortstack{специ-\\фична} & \shortstack{въз-\\можна} & силен релационен договор с доказателства за изпълнение на живо \\
MySQL & налична & налична & налична & \shortstack{специ-\\фична} & \shortstack{въз-\\можна} & експлоатационният път е наличен, но детайлите за свързване и диалекта остават специфични за конектора \\
MongoDB & липсва & липсва & липсва & аналог & \shortstack{не е\\централна} & силен документен модел без профил на възможностите в SQL стил \\
Redis & липсва & липсва & липсва & аналог & \shortstack{ограни-\\чена} & семантика ключ-стойност с минималистичен, но честен договорен профил \\
Neo4j & аналог & налична & липсва & \shortstack{не е\\централна} & \shortstack{не е\\централна} & графова целева система с фокус върху транзакциите, а не върху слой за SQL съединяване/агрегация \\
Cassandra & липсва & налична & липсва & \shortstack{не е\\централна} & \shortstack{въз-\\можна} & широка колонна система с ограничен, но формално дефиниран набор от възможности \\
\bottomrule
\end{tabularx}
\end{table}

Таблица~\ref{tab:backend_capability_matrix} позволява по-прецизно да се види къде библиотеката е най-близо до класическо ОРС поведение и къде съзнателно работи с частичен договор. Релационните целеви системи концентрират по-богат профил на възможностите, докато документните, графовите и ключ-стойност вариантите се оценяват през други критерии. Именно тази контролирана нееднородност прави дизайна за множество свързващи компоненти инженерно честен.

\subsection{Съпоставка между тестови доказателства и договор на компонента}

\begin{table}[H]
\centering
\small
\caption{Свързващи компоненти, профил на възможностите и налична тестова опора}
\label{tab:backend_evidence}
\begin{tabularx}{\textwidth}{L{2.1cm}Q{2.4cm}L{2.8cm}Y}
\toprule
\textbf{Компонент} & \textbf{Профил на възможностите} & \textbf{Модулен тест/тест на договора} & \textbf{Интеграционен/жизнен тест} \\
\midrule
SQLite & съединявания, транзакции, агрегация & утвърждавания на възможностите в \path{tests/интеграционни/test_sqlite.cpp} & \path{tests/интеграционни/test_sqlite.cpp} \\
PostgreSQL & съединявания, транзакции, агрегация & \path{tests/модулни/test_postgresql_конектор.cpp} & \path{tests/интеграционни/test_postgresql_live.cpp} \\
MySQL & съединявания, транзакции, агрегация & \path{tests/модулни/test_mysql_конектор.cpp} & \path{tests/интеграционни/test_mysql_live.cpp} \\
MongoDB & без labelи за SQL възможности & \path{tests/модулни/test_mongodb_конектор.cpp} & \path{tests/интеграционни/test_mongodb_live.cpp} \\
Redis & без съединявания/транзакции/агрегация & \path{tests/модулни/test_redis_конектор.cpp} & \path{tests/интеграционни/test_redis_live.cpp} \\
Neo4j & транзакции & \path{tests/модулни/test_neo4j_конектор.cpp} & \path{tests/интеграционни/test_neo4j_live.cpp} \\
Cassandra & транзакции & \path{tests/модулни/test_cassandra_конектор.cpp} & \path{tests/интеграционни/test_cassandra_live.cpp} \\
\bottomrule
\end{tabularx}
\end{table}

Таблица~\ref{tab:backend_evidence} е важна, защото превежда сравнението от равнището на описателен анализ (описателен анализ) към равнището на проверими доказателства от хранилището (проверими доказателства от хранилището). За всеки компонент има поне един ясен корелат на ниво договор или на ниво изпълнение на живо. Това намалява риска главата да се превърне в декларативен преглед без техническа опора.

\subsection{Сравнителен синтез}

Съпоставката между целевите системи показва, че библиотеката успява да запази единен логически модел на равнището на типове същности, полета, заместители и междинно представяне. Различията се появяват на равнището на възможностите и на физическата организация на данните. Това е и най-важният извод на настоящата глава: успешната ОРС архитектура за множество компоненти не премахва различията между моделите за съхранение, а ги прави контролируеми чрез формален договор на конектора.

\begin{table}[H]
\centering
\small
\caption{Съпоставка на материализацията в целевите системи}
\label{tab:backend_materialisation}
\begin{tabularx}{\textwidth}{L{2.0cm}L{2.5cm}L{2.8cm}Y}
\toprule
\textbf{Компонент} & \textbf{Основна структура} & \textbf{Логическа интерпретация} & \textbf{Показателно ограничение} \\
\midrule
SQLite & таблица & пълна SQL ОРС интерпретация & стандартна релационна дисциплина и богата изразителност на заявките \\
PostgreSQL/MySQL & таблица & преносимо SQL представяне & различен модел на свързване и диалектни особености \\
MongoDB & документ & интерпретация като филтър/проекция & частична еквивалентност спрямо SQL операции \\
Redis & ключ/стойност, хеш (hash) & базиран на ключове достъп до същности & търсене предимно по първичен ключ и липса на семантика за съединяване \\
Neo4j & възел/свойства & интерпретация като графов модел (graph-pattern) & специфична за графове семантика на заявките вместо слой за SQL съединяване \\
Cassandra & широк колонен ред & CQL с ключови ограничения & достъп, управляван от разделянето, и ограничен профил на предикатите \\
\bottomrule
\end{tabularx}
\end{table}

В резултат дизайнът на библиотеката за множество свързващи компоненти може да бъде защитен с по-прецизна формулировка. Системата не обещава ``универсална база данни през един програмен интерфейс'', а \emph{унифициран логически слой с експлицитни граници на целевата система}. Именно тази формулировка е едновременно инженерно честна и научно обоснована.
